\section{Agent System}
\label{sec:agent_system}

The agent system is the core reasoning engine of \name, implementing an extended ReAct-based~\cite{yao2023react} loop that orchestrates multi-model routing, multi-phase reasoning, tool execution, and context management. This section explores the design space for each component, explaining why particular approaches were chosen, what alternatives were considered, and what lessons were learned during development.

\subsection{Workload-Optimized Multi-Model Architecture}
\label{sec:multi_model}

Different execution phases benefit from different model capabilities. Reasoning tasks benefit from extended thinking without tool distraction. Visual tasks require vision-language models. Bulk summarization benefits from cheaper, faster models. Using a single model for all tasks either wastes cost (using expensive models for simple tasks) or sacrifices quality (using cheap models for complex reasoning). Alternatives include using a single model for everything (simple but inflexible) and ensemble execution (maximum quality but prohibitive latency and cost). \name adopts task-specific routing, directing different execution phases to specialized models at the cost of selection logic complexity.

\paragraph{Five model roles with fallback chains.} Five distinct workload categories route to specialized models:

\begin{packeditemize}
\item \textbf{Action model:} Primary execution model for tool-based reasoning. Default for all workloads unless specialized model specified.

\item \textbf{Thinking model:} Optional model for extended reasoning without tool access. Enables focus on strategic planning without tool-calling pressure. Fallback: action model.

\item \textbf{Critique model:} Optional model for self-evaluation. Inspired by Reflexion~\cite{shinn2023reflexion} but applied selectively rather than on every turn. Fallback: thinking model $\to$ action model.

\item \textbf{Vision model:} Vision-language model for processing screenshots and images. Essential for visual debugging tasks. Fallback: action model if vision-capable.

\item \textbf{Compact model:} Smaller, faster model for summarization during context compaction. Prioritizes speed and cost over reasoning depth. Fallback: action model.
\end{packeditemize}

\paragraph{Provider abstraction and caching.} Each model selection triggers lazy initialization of provider-specific API clients. This reduces startup latency - only models actually used in a session are initialized. Model capabilities (context length, vision support, reasoning features) are cached locally with time-to-live refresh, enabling offline startup and background updates following a stale-while-revalidate pattern.

Capabilities are cached because cold startup without cache requires API calls, creating a network dependency. The cache trades potential staleness (model specifications may change between refreshes) for startup reliability (the system works offline), with background refresh ensuring eventual consistency.

\paragraph{Design evolution.} Eager initialization of all models at startup caused slow launches and wasted memory; lazy initialization solved this. Early versions had separate code paths for different providers, leading to frequent API incompatibilities; a unified adapter interface eliminated this. Network dependency on every launch was fragile; caching with background refresh provided both reliability and freshness.

\subsection{Extended Four-Phase Reasoning Pipeline}
\label{sec:react_loop}

\Cref{sec:react_executor} presented the execution pipeline from an implementation perspective. This subsection analyzes the design space and reasoning behind each phase, and provides the algorithmic pseudocode.

Standard ReAct~\cite{yao2023react} interleaves reasoning and action, where the model thinks and acts in the same turn. This conflation limits deliberation: tool schemas consume context, creating pressure to act quickly rather than think deeply. Moreover, finite context windows eventually overflow in long conversations. Alternatives include pure ReAct (simple but biased toward premature action), chain-of-thought prompting (inflexible across task complexity), and Reflexion-inspired self-critique loops (maximum quality but doubled latency). \name combines a separate thinking phase with optional self-critique: explicit tool-free reasoning before action enables depth control and prevents premature tool use.

\paragraph{Four-phase pipeline.} The system extends ReAct with explicit phases for compaction, thinking, and self-critique. \Cref{alg:react} presents the complete loop structure.

\begin{algorithm}[t]
\caption{Extended Four-Phase ReAct Loop}
\label{alg:react}
\begin{algorithmic}[1]
\Require User message $m$, Agent $\mathcal{A}$, Tool registry $\mathcal{T}$, Session $\mathcal{S}$
\Ensure Response summary, error status, latency
\State $\mathcal{S}.\text{add}(m)$ \Comment{Persist user message}
\State $\text{nudge\_count} \gets 0$
\Repeat
    \State \textbf{// Phase 0: Auto-Compaction}
    \If{$\text{token\_count}(\mathcal{S}) > 0.7 \times \text{max\_context}$}
        \State $\mathcal{S} \gets \text{compact}(\mathcal{S})$ \Comment{Summarize middle, preserve recent}
    \EndIf
    \State \textbf{// Phase 1: Thinking (if enabled)}
    \If{$\text{thinking\_level} \neq \text{OFF}$}
        \State $\text{trace} \gets \mathcal{A}.\text{call\_thinking\_llm}(\mathcal{S})$ \Comment{No tools}
        \State $\mathcal{S}.\text{add\_trace}(\text{trace})$
        \If{$\text{thinking\_level} = \text{SELF\_CRITIQUE}$}
            \State \textbf{// Phase 2: Self-Critique}
            \State $\text{critique} \gets \mathcal{A}.\text{call\_critique\_llm}(\text{trace})$
            \State $\text{trace} \gets \mathcal{A}.\text{refine\_with\_critique}(\text{trace}, \text{critique})$
        \EndIf
    \EndIf
    \State \textbf{// Phase 3: Action}
    \State $\text{response}, \text{tool\_calls} \gets \mathcal{A}.\text{call\_llm}(\mathcal{S}, \mathcal{T})$ \Comment{With tools}
    \If{$\text{tool\_calls} \neq \emptyset$}
        \For{$\text{tc} \in \text{tool\_calls}$}
            \State $\text{result} \gets \mathcal{T}.\text{execute}(\text{tc})$
            \State $\mathcal{S}.\text{add}(\text{tc}, \text{result})$
        \EndFor
    \Else
        \If{last tool failed $\land$ nudge\_count $< 3$}
            \State $\text{nudge} \gets \text{get\_smart\_nudge}(\text{error})$ \Comment{Context-aware recovery}
            \State $\mathcal{S}.\text{add}(\text{nudge})$
            \State $\text{nudge\_count} \gets \text{nudge\_count} + 1$
        \Else
            \State \textbf{break} \Comment{Implicit completion}
        \EndIf
    \EndIf
\Until{task\_complete called $\lor$ max iterations reached}
\State \Return summary, error, latency
\end{algorithmic}
\end{algorithm}

\paragraph{Phase 0: Auto-Compaction.}
Before each iteration, the loop checks whether the accumulated messages exceed a configurable fraction of the context window. If so, it compresses history by preserving the system prompt and recent messages while summarizing the middle. The full compaction mechanism, including the adaptive observation lifecycle and threshold selection rationale, is described in \Cref{sec:adaptive_compaction}.

\paragraph{Phase 1: Thinking.}
The thinking phase calls a dedicated model without tool schemas, allowing extended reasoning free from tool-calling pressure. Separating thinking from action prevents premature tool use: when tools are available, models tend to act quickly rather than think deeply. Six configurable depth levels - from OFF to DEEP - let users balance latency against deliberation quality on a per-task basis.

\paragraph{Phase 2: Self-Critique (Optional).}
Inspired by Reflexion~\cite{shinn2023reflexion}, this optional phase evaluates the thinking trace for logical soundness and identifies risks before the action phase begins. A critique model reviews the trace, then the thinking model refines its reasoning with the critique as additional input. Because critique doubles thinking latency, it is activated selectively - justified for complex tasks with high failure cost but skipped for simple operations.

\paragraph{Phase 3: Action.}
The action phase follows standard ReAct: the primary model receives the complete message history (including any thinking trace) plus all available tool schemas and produces tool calls or a text response. Tools execute sequentially by default, with parallel execution supported for independent operations such as multiple file reads or subagent spawning. Completion is detected through three mechanisms: explicit (the agent calls a completion tool), implicit (a text response with no tool calls after successful execution), or forced (iteration or nudge limit reached).

\paragraph{Design evolution.} The original design had no automatic compression, leading to frequent context overflows. We considered adaptive thresholds based on message size trends but found a fixed threshold sufficient and simpler. Early experiments applied self-critique on every thinking output; this proved too slow for routine operations, motivating selective activation. Similarly, heuristic compression (dropping tool results, keeping user messages) was replaced by LLM-based summarization, which better preserves semantic information at higher compression ratios.

\subsection{Dual-Agent Safety Architecture}
\label{sec:dual_agent}

\Cref{sec:agent_core} describes the dual-mode architecture and its operational flow. Here we note the key design insight: safety is enforced at the schema level rather than through runtime permission checks. The planning agent's tool schema contains only read-only tools; write operations are not filtered but absent. This eliminates an entire class of bypass concerns - the agent cannot attempt writes because the tools do not exist in its schema. Conversation history is preserved across mode transitions, enabling a seamless ``plan, review, execute'' workflow.

\paragraph{Design evolution.} The initial design used a single agent with runtime permission checks. We observed bypass attempts where the agent tried write operations during planning, and the filtering logic was complex and error-prone. Separate agent instances with distinct tool schemas made enforcement automatic. We also found that the planning agent needed a bounded iteration budget - without one, exploration could run unnecessarily long.

\subsection{Subagent Orchestration}
\label{sec:subagents}

Some tasks benefit from focused expertise (codebase exploration, user clarification) while others require coordination with main agent state (task management). The question is when specialized agents should handle subtasks versus having the main agent do everything.

\paragraph{Design pattern: Hierarchical delegation.} The main agent can spawn specialized subagents for specific subtasks, each with filtered tool access and specialized prompts. Subagents execute in isolated contexts with their own iteration budgets, preventing unbounded execution.

\paragraph{Subagent specialization.} Different subagents serve different roles:

\begin{packeditemize}
\item \textbf{Code explorer:} Read-only tools for codebase navigation. Specialized for understanding existing code structure.
\item \textbf{Strategic planner:} Read-only tools plus extended reasoning. Focuses on high-level planning without execution.
\item \textbf{Web tools:} Web fetching and file writing for cloning web content. Combines retrieval with persistence.
\item \textbf{User clarification:} Minimal toolset for gathering input. Focuses on eliciting information without distraction.
\end{packeditemize}

\paragraph{Tool filtering rationale.} Subagent capabilities are deliberately restricted for three reasons. First, task management tools are excluded from subagents so that only the main agent coordinates todo lists, preventing race conditions and inconsistent state. Second, limited tool access reduces context size and focuses each subagent on its specific role - an exploration subagent does not need write capabilities. Third, restricted tools limit the blast radius of errors: an exploration subagent cannot accidentally modify files.

\paragraph{Parallel execution.} The main agent can spawn multiple subagents concurrently (each in its own thread) for independent queries such as parallel file searches, codebase exploration, or web fetches. This trades thread overhead for latency reduction on independent operations.

\paragraph{Completion detection.} Subagent signals completion through marker in output. Main agent receives subagent results and continues with its own reasoning.

\paragraph{Design evolution.} Early versions gave subagents the same tools as the main agent. This caused context pollution, role confusion, and conflicts when both agents attempted to update todos simultaneously. Restricting each subagent's tool set to its specific role improved both focus and efficiency.

\subsection{Template-Based Error Recovery}
\label{sec:error_recovery}

Guiding agents toward recovery without hard-coded logic is difficult. Hard-coded retry rules quickly become combinatorially complex (different recovery strategies for different error types in different contexts). Generic error messages lack actionable guidance. Alternatives include explicit if/else for each error type (inflexible as error types grow) and returning raw errors to the agent (uninformative). \name adopts template-based nudges: classify the error, retrieve a targeted recovery template, and inject it as a system message, centralizing recovery knowledge while remaining extensible.

\paragraph{Smart nudge mechanism.} When a tool fails and the agent attempts completion without retry, the system applies a four-step recovery process:

\begin{enumerate}
\item \textbf{Classify error:} Pattern-match error message to category (permission error, file not found, edit mismatch, syntax error, rate limit, timeout).
\item \textbf{Retrieve template:} Fetch context-aware recovery guidance from template repository.
\item \textbf{Format with context:} Inject context-specific details (file paths, error messages) into template placeholders.
\item \textbf{Apply as system message:} Provide targeted recovery strategy without hard-coded logic.
\end{enumerate}

A budget of 3 nudge attempts provides multiple recovery chances while preventing infinite retry loops - after three failed attempts, the agent accepts the failure gracefully rather than looping indefinitely. Templates centralize recovery knowledge in one place, making strategies extensible (new error types require no code changes) and customizable (users can modify templates for project-specific guidance).

\paragraph{Design evolution.} The initial design used direct retry logic for each error type, which suffered combinatorial explosion as error types multiplied. Consolidating all strategies into templates made recovery easier to understand, modify, and extend. Early versions also had no nudge budget, which led to infinite loops when the agent could not recover; the budget forces eventual progress.

\subsection{Context-Aware System Reminders}
\label{sec:system_reminders}

Error recovery nudges (\Cref{sec:error_recovery}) are one instance of a broader mechanism we call \emph{system reminders}: short, targeted messages injected into the conversation as \texttt{role:~user} messages at specific decision points within the ReAct loop. This subsection presents the full architecture - its motivation, design, and the six categories of reminders that keep the agent on track across long-horizon tasks.

\paragraph{Attention decays over distance.} LLMs process the system prompt at position zero of the context window. As the conversation grows - tool calls accumulate, file contents are read, code is written - the system prompt drifts further from the model's active attention. The instructions are still present in the context, but their influence on generation weakens with distance: a well-documented property of transformer attention where tokens attend more strongly to nearby positions than to distant ones. The consequence is predictable and reproducible: an agent instructed to ``always run tests after editing'' stops running tests after 15+ tool calls; an agent instructed to ``complete all todos before finishing'' declares victory with half the list incomplete; an agent that encounters a tool failure gives up rather than retrying with corrected input. The system prompt said the right thing. The agent stopped listening.

\paragraph{Design space.}

\begin{packeditemize}
\item \textbf{Front-loaded prompting:} Pack every rule into the system prompt. By turn 40, the relevant constraint has decayed below the effective attention threshold. The agent violates the rule.

\item \textbf{Periodic re-injection:} Re-insert the full system prompt every $N$ turns. Wastes tokens on instructions irrelevant to the current moment and does not guarantee proximity to the decision point.

\item \textbf{Event-driven just-in-time injection (adopted):} Inject a short, single-concern reminder immediately before the decision point where the agent needs it. Maximum recency, minimal token cost, context-specific content.
\end{packeditemize}

\paragraph{Three-tier context architecture.} Guidance is organized into three tiers, each operating at a different timescale (\Cref{fig:system_reminders} illustrates concrete examples from Tiers~1 and~2):

\begin{packedenumerate}
\item \textbf{Tier~1 - Static system prompt.} Assembled once per turn by \texttt{PromptComposer} from 18 modular markdown sections (\Cref{sec:prompt_composition}). Defines \emph{who} the agent is: identity, capabilities, security policy, tool-use rules. Delivered as \texttt{role:~system}.

\item \textbf{Tier~2 - Dynamic system reminders.} Injected at decision points within the ReAct loop as \texttt{role:~user} messages. Addresses \emph{what to do right now}: nudges, signals, corrections. 24 named reminders organized into six categories (described below).

\item \textbf{Tier~3 - Long-horizon persistence.} Operates across turns and sessions: automatic context compaction (\Cref{sec:compaction}), project-scoped session persistence (\Cref{sec:session_management}), dual-memory architecture (\Cref{sec:dual_memory}), todo enforcement, and validated message lists. Preserves \emph{what happened before} beyond the context window.
\end{packedenumerate}

Tier~1 establishes the rules. Tier~2 enforces them at the moments they matter. Tier~3 ensures context survives across compaction boundaries and session restarts.

\begin{figure}[htbp]
    \centering
    \begin{subfigure}[b]{0.95\linewidth}
        \includegraphics[width=\linewidth]{figures/system_reminder_1.pdf}
        \caption{A \emph{task lifecycle signal}: after the user approves the plan, the reminder restates the approved plan content, lists the created todos, and provides explicit workflow instructions (``\emph{mark todo as doing, implement, mark as done}''). Without this reminder, the agent acknowledges the approval and waits instead of beginning implementation.}
        \label{fig:reminder_1}
    \end{subfigure}

    \vspace{1em}
    \begin{subfigure}[b]{0.95\linewidth}
        \includegraphics[width=\linewidth]{figures/system_reminder_2.pdf}
        \caption{An \emph{error recovery nudge}: after \texttt{edit\_file} fails with ``\emph{old\_content not found},'' the agent produces text with no tool calls (giving up). The loop classifies the error as \texttt{edit\_mismatch} and injects specific recovery guidance: ``\emph{read the file again to get exact current content, then retry.}'' Without this nudge, the agent apologizes and abandons the edit.}
        \label{fig:reminder_2}
    \end{subfigure}
    \caption{Examples of context-aware system reminders injected during the reasoning loop. Each reminder is a short \texttt{role:~user} message placed at maximum recency - immediately before the next LLM call - to counteract attention decay from the distant system prompt.}
    \label{fig:system_reminders}
\end{figure}

\paragraph{Rationale for \texttt{role:~user}.} System messages establish identity at conversation start. After 40 turns, another system message gets lost because the model has already internalized (and partially forgotten) its system prompt. A \texttt{role:~user} message appears in the flow of dialogue at the position of highest recency; the model treats it as something that just happened, something that demands a response. The reminders do not add new rules; they nudge the agent at the moment it needs guidance.

\paragraph{Storage and retrieval.} All 24 reminder templates live in a single file (\texttt{reminders.md}) using \texttt{-{}-{}-~section\_name~-{}-{}-} delimiters to define named sections. A single entry point, \texttt{get\_reminder(name, **kwargs)}, parses the file once into a module-level cache, looks up the section by name, and runs \texttt{str.format(**kwargs)} for placeholder substitution (e.g., \texttt{\{count\}}, \texttt{\{todo\_list\}}, \texttt{\{thinking\_trace\}}). Standalone \texttt{.txt} files serve as fallbacks for longer templates. This design centralizes all reminder text in one human-readable file, making it easy to audit, edit, and extend.

\paragraph{Reminder catalog: six categories.} The 24 named reminders serve distinct purposes:

\begin{packeditemize}
\item \textbf{Phase control} (4 reminders): Manage thinking/action phase transitions. Inject thinking traces into the action phase context; control whether the agent should reason deeply or act directly.

\item \textbf{Task lifecycle} (5 reminders): Steer phase transitions in multi-step workflows. After a subagent returns, nudge the main agent to synthesize results and continue (\Cref{fig:reminder_1} shows the plan-approval variant). After plan approval, restate the plan and todos with explicit workflow instructions. On session resume, reference the existing plan file.

\item \textbf{Todo enforcement} (2 reminders): Gate premature completion. When the agent calls \texttt{task\_complete} with incomplete todos, reject the call and list remaining items. When all todos are done, signal the agent to finalize. Capped at 2 nudges per run - after two attempts, accept the agent's judgment.

\item \textbf{Error recovery} (8 reminders): The mechanism described in \Cref{sec:error_recovery}. A generic nudge plus six type-specific nudges selected by error classification (permission, file-not-found, syntax, rate-limit, timeout, edit-mismatch) and one Docker-specific nudge (\Cref{fig:reminder_2} shows the edit-mismatch variant).

\item \textbf{Behavioral} (5 reminders): Correct observable anti-patterns. After 5+ consecutive read-only tool calls, break the exploration spiral. After the user denies a tool call, prevent re-attempt of the same command. On empty completion (no text, no tool calls), request a brief outcome summary. At the iteration safety limit, force summarization.

\item \textbf{JSON retry} (2 reminders): Specialized parse-retry prompts for the memory system's Reflector and Curator components when JSON output fails to parse.
\end{packeditemize}

\paragraph{Injection timing within the ReAct loop.} Reminders injected by the executor follow a strict 9-step ordering within each iteration of the loop (\Cref{alg:react}): (1)~auto-compaction check, (2)~interrupt check, (3)~thinking phase with optional trace injection, (4)~subagent-completion signal, (5)~drain UI-thread messages, (6)~interrupt check, (7)~action phase LLM call, (8)~response dispatch - where reminders diverge based on whether the response contains tool calls or not, and (9)~session persistence. On the no-tool-calls path, the system checks for failed-tool nudges, incomplete-todo nudges, and empty-completion nudges, in that order. On the has-tool-calls path, it checks for plan-approved signals, all-todos-complete signals, tool-denied nudges, and consecutive-reads nudges after tool execution.

\paragraph{Safety guards.} Unguarded reminders degenerate into noise. Two guard mechanisms prevent this:

\begin{packeditemize}
\item \textbf{One-shot flags:} Boolean flags in the iteration context ensure certain reminders fire at most once per agent run: \texttt{plan\_approved\_signal\_injected}, \texttt{all\_todos\_complete\_nudged}, \texttt{completion\_nudge\_sent}. A nudge that fires every iteration stops being a nudge and becomes background static.

\item \textbf{Attempt budgets:} Reminders that may need to fire multiple times are capped: \texttt{MAX\_TODO\_NUDGES~=~2} (after two attempts to address incomplete todos, accept the agent's judgment), \texttt{MAX\_NUDGE\_ATTEMPTS~=~3} (after three recovery attempts for a failed tool, accept the failure gracefully).
\end{packeditemize}

\paragraph{Graceful degradation.} If a reminder template is missing or retrieval fails, the agent still has the system prompt from Tier~1. Reminders are a safety net, not the primary guidance mechanism - the system works without them; it works better with them.

\paragraph{Design evolution.} The initial system relied entirely on the system prompt. In long sessions (30+ tool calls), agents reliably exhibited attention-decay failures: premature completion, exploration loops, and failure to recover from errors. Adding just-in-time reminders resolved each failure mode. Early error recovery used a single generic ``try again'' message; classifying errors into six types with specific guidance (\Cref{sec:error_recovery}) substantially improved recovery rates - ``read the file again'' is more actionable than ``fix the issue.'' Before one-shot flags and attempt budgets, some reminders fired on every iteration, causing the agent to loop on the nudge itself; guards were essential for stability. Early experiments also injected reminders as \texttt{role:~system} messages, which were less effective because the model treated them as background instructions rather than conversational prompts demanding a response.

\subsection{Completion Detection}
\label{sec:completion}

\paragraph{Completion mechanisms.} Three mechanisms:

\begin{packedenumerate}
\item \textbf{Explicit completion:} Agent invokes completion tool with summary and status (success/failure). Clearest signal of intentional completion.
\item \textbf{Implicit completion:} Agent produces text response with no tool calls after successful tool execution. Interprets conversational response as task completion.
\item \textbf{Forced completion:} Nudge budget exhausted or iteration limit reached. Prevents infinite execution.
\end{packedenumerate}

\paragraph{Task validation before completion.} If outstanding task items remain, agent receives additional nudges to address them before forced completion. Prevents premature completion with incomplete work.

Validation was added because early versions allowed premature completion with tasks left unfinished. The validation step ensures the agent either completes outstanding work or explicitly acknowledges incomplete tasks before terminating.

\subsection{Modular System Prompt Composition}
\label{sec:prompt_composition}

\Cref{sec:initialization} describes the core composition mechanism: priority-ordered conditional sections assembled at runtime through a filter--sort--load--join pipeline (\Cref{fig:prompt_composer_init}). This subsection covers the additional design aspects not addressed there.

\paragraph{Mode-specific variants.} Different reasoning phases require fundamentally different prompts. Normal (action) mode loads all 18 sections across four tiers. Thinking mode, a tool-free reasoning pre-phase (\Cref{sec:react_loop}), loads only 4 purpose-built sections: available-tool awareness (so the model knows what actions are possible without being tempted to invoke them), subagent guidance, code-reference conventions, and output-format rules. No tool-selection or code-quality sections are included, because the thinking phase produces reasoning traces rather than tool calls. Planning mode uses a standalone template optimized for read-only exploration. This mode specialization is achieved through a factory function: \texttt{create\_composer(templates\_dir, mode="system/main")} returns the 18-section action composer, while \texttt{mode="system/thinking"} returns the 4-section thinking composer.

\paragraph{Variable substitution.} Templates use \texttt{\$\{VAR\}} placeholders resolved at render time by a \texttt{PromptRenderer}. A centralized \texttt{PromptVariables} registry maps symbolic names to concrete tool identifiers - for example, \texttt{\$\{EDIT\_TOOL.name\}} resolves to \texttt{edit\_file}. This indirection decouples template prose from tool naming: renaming a tool requires changing one registry entry rather than editing every template. The renderer also supports dotted attribute access and runtime variables for session-specific values.

\paragraph{Two-tier fallback.} If an individual section file is missing, the composer skips it and proceeds - the agent starts with a slightly reduced prompt rather than failing. If modular composition fails wholesale (e.g., the templates directory is absent), the builder falls back to a monolithic core template. This guarantees agent startup under partial-deployment conditions.

\paragraph{Design evolution.} The initial design used a single large template where conditional blocks via string interpolation led to fragile nested logic. Including git-workflow instructions in non-git projects wasted hundreds of tokens with no benefit; conditional predicates eliminated this waste. Early thinking-mode prompts included tool-use sections, biasing the model toward premature tool calls; maintaining separate section registries per mode cleanly eliminated this interference.
